% SEGMENT OLP-0634-S01
% Part: history
% Chapter: set-theory
% Section: infinitesimals

\documentclass[../../../include/open-logic-section]{subfiles}

\begin{document}

\olfileid{his}{set}{infinitesimals}
\olsection{நுண்ணளவுகளும் வகையிடலும்}

நியூட்டனும் லைப்னிட்ஸும் பதினேழாம் நூற்றாண்டின் இறுதியில், தனித்தனியாக,
நுண்கணிதத்தைக் கண்டறிந்தனர். நுண்கணிதத்தின் முக்கியமான பயன்பாடுகளில் ஒன்று
\emph{வகையிடல்}. பொதுவாகச் சொன்னால், ஒரு சார்பின் குறிப்பிட்ட புள்ளியில்
அதன் “மாறுபாட்டு வீதம்” அல்லது சாய்வு என்ன என்பதை வகையிடல்
விளக்குகிறது.

% SEGMENT OLP-0634-S02
இந்தக் கருத்தைத் தெளிவாகக் காட்ட ஒரு வழி இருக்கிறது. கீழே கருப்பு நிறத்தில்
வரையப்பட்டுள்ள $f(x) = \nicefrac{x^2}{4} + \nicefrac{1}{2}$ என்ற
சார்பைக் கவனியுங்கள்:
\begin{center}
	\begin{tikzpicture}[scale=1]
	\draw[->, gray] (-1,0) -- (4.25,0) node[right] {$x$};
	\draw[->, gray] (0,-0.25) -- (0,5.25) node[above] {$f(x)$};

	\foreach \x/\xtext in {1/1, 2/2, 3/3, 4/4}
	\draw[shift={(\x,0)}, gray] (0pt,2pt) -- (0pt,-2pt) node[below] {$\xtext$};

	\foreach \y/\ytext in {1/1, 2/2, 3/3, 4/4, 5/5}
	\draw[shift={(0,\y)}, gray] (2pt,0pt) -- (-2pt,0pt) node[left] {$\ytext$};

	\draw[oldiagcolorC] (0.5, 9/16) -- (3.5, 57/16) -- (3.5, 9/16)--cycle;
	\draw[oldiagcolorD] (.5, 9/16) -- (2.5, 33/16) -- (2.5, 9/16) -- cycle;
	\draw[oldiagcolorE] (.5, 9/16) -- (1.5, 17/16) -- (1.5, 9/16) -- cycle;
	\draw[thick] (-1,.75) parabola bend (0,0.5) (4,4.5);
	\end{tikzpicture}
\end{center}

% SEGMENT OLP-0634-S03
$c = \nicefrac{1}{2}$ என்ற புள்ளியில் சார்பின் சாய்வைக் கண்டுபிடிக்க
விரும்புகிறோம் என்று வைத்துக்கொள்வோம். முதலில், அந்தப் புள்ளியில் உள்ள
சாய்வைத் தோராயமாகக் காட்டும் கர்ணத்தையுடைய ஒரு முக்கோணத்தை வரைகிறோம்;
மேலுள்ள சிவப்பு முக்கோணம் அதற்கு ஓர் எடுத்துக்காட்டு. நமது
முக்கோணத்தின் அடிப்பக்க நீளம் $\beta$ எனில், அதன் உயரம்
$f(\nicefrac{1}{2}+\beta) - f(\nicefrac{1}{2})$. எனவே அதன்
கர்ணத்தின் சாய்வு:
\[
\frac{f(\nicefrac{1}{2}+\beta) - f(\nicefrac{1}{2})}{\beta}.
\]
ஆகவே, அடிப்பக்க நீளம்~$3$ கொண்ட !!{colorC} முக்கோணத்தின் சாய்வு
துல்லியமாக~$1$. அடிப்பக்க நீளம்~$2$ கொண்ட சிறிய !!{colorD}
முக்கோணத்தின் கர்ணம் இன்னும் நல்ல தோராயத்தைத் தருகிறது; அதன் சாய்வு
$\nicefrac{3}{4}$. அடிப்பக்க நீளம்~$1$ கொண்ட, அதைவிடச் சிறிய
!!{colorE} முக்கோணம் மேலும் நல்ல தோராயத்தைத் தருகிறது; அதன் சாய்வு
$\nicefrac{1}{2}$.

% SEGMENT OLP-0634-S04
முக்கோணங்கள் சிறிதாகச் சிறிதாக ஆகும்போது தோராயங்கள் மேம்படுகின்றன. எனவே
இவ்வாறு சொல்லத் தோன்றலாம்: \emph{நுண்ணளவான} அடிப்பக்க நீளமுள்ள
முக்கோணத்தின் கர்ணமே $c = \nicefrac{1}{2}$ புள்ளியில் உள்ள சாய்வைத்
தரும். இதன் மூலம் $c$ புள்ளியில் $f$ சார்பின் (முதல்) வகைக்கெழுவுக்கான
சூத்திரம் கிடைப்பதாகக் கூறலாம்:
\[
{f'}(c) = \frac{f(c+\beta) - f(c)}{\beta} \text{ இங்கு $\beta$ நுண்ணளவானது.}
\]
தோராயமாகச் சொன்னால், நியூட்டனும் லைப்னிட்ஸும் கூறியது இதைத்தான்.

ஆனால் அவர்கள் அப்படிக் கூறியதும், \emph{நுண்ணளவு} என்றால் என்ன என்று
கேட்க வேண்டியிருக்கிறது. இங்கே ஒரு தீவிரமான இருமுனைச் சிக்கல்
எழுகிறது. $\beta = 0$ எனில், $0$-ஆல் வகுக்க வேண்டி வருவதால் $f'$
வரையறுக்கப்படாது. மாறாக $\beta >
0$ எனில், சாய்வின் \emph{தோராயம்தான்} கிடைக்கும்; சாய்வே கிடைக்காது.

% SEGMENT OLP-0634-S05
இது பிற்காலத்திலிருந்து பின்னோக்கிப் புகுத்தப்பட்ட கவலை அல்ல.
நியூட்டனின் வழிவந்தவர்களை விமர்சித்த பெர்க்லி கூறியதைக் கவனியுங்கள்:
\begin{quote}
குறிகள் ஒன்றையோ ஒன்றுமில்லாததையோ குறிக்கலாம் என்பதை நான் ஒப்புக்கொள்கிறேன்.
அதன்படி, முதலில் எழுதப்பட்ட $c + \beta$ என்ற குறியீட்டில் $\beta$
என்பது ஓர் அதிகரிப்பையோ ஒன்றுமில்லாததையோ குறித்திருக்கலாம். ஆனால் அது
இவற்றில் எதைக் குறிக்கிறது என்று நீங்கள் தேர்ந்தெடுத்தாலும், அந்தப்
பொருளோடு முரணின்றி வாதிட வேண்டும்; ஒரே குறிக்கு இரு பொருள்களை மாறி மாறிக்
கொண்டு வாதத்தைத் தொடரக் கூடாது. அப்படிச் செய்வது வெளிப்படையான
குதர்க்கமாகும். (\citealt[\S{}XIII]{Berkeley1734};
முன்னுள்ள உரையுடன் பொருந்துமாறு மாறிகள் மாற்றப்பட்டுள்ளன)
\end{quote}
பெர்க்லியின் விமர்சனத்திலிருந்து நுண்ணளவுக் கணிதத்தைக் காக்க,
“நுண்ணளவுகள்” என்று சொல்வது வெறும் உருவகமே என்று ஒருவர் பதிலளிக்கலாம்.
\emph{மிகவும் சிறிய} முக்கோணத்தை எடுத்தால், தொடுகோட்டுக்குப்
\emph{போதுமான அளவு நல்ல} தோராயம் கிடைக்கும் என்று கூறலாம். அதற்கும்
பெர்க்லி பதில் அளித்தார்: பொறியியலுக்கு அது போதுமானதாக இருந்தாலும்,
கணிதத்தின் \emph{நிலையை} அது குலைக்கிறது; ஏனெனில்
\begin{quote}
\emph{in rebus mathematicis errores qu\`{a}m minimi
non sunt contemnendi} என்று நமக்குச் சொல்லப்படுகிறது. [கணிதத்தில்,
மிகச் சிறிய பிழைகளைக்கூடப் புறக்கணிக்கக் கூடாது.]
\citep[\S{}IX]{Berkeley1734}
\end{quote}
சாய்வெழுத்தில் உள்ள இலத்தீன் பகுதி, நியூட்டன் 1704-இல் எழுதிய
\emph{Quadrature of Curves} நூலிலிருந்து கிட்டத்தட்ட சொற்சொல்லாக
எடுக்கப்பட்ட மேற்கோள்.

பெர்க்லியின் மெய்யியல் ஆட்சேபங்கள் ஆழமும் கூர்மையும் கொண்டவை.
இருந்தாலும், நுண்கணிதம் மிகப் பெரும் வெற்றி பெற்ற முயற்சியாக அமைந்தது;
கணிதவியலாளர்கள் பிழையில் சிக்காமல் தொடர்ந்து அதைப் பயன்படுத்தினர்.

\end{document}
